\section{Introduction and Proof Outline}
\label{sec:introduction}

\subsection{The geometric obstruction}

We prove that seven open equilateral triangles of side one cannot cover a
regular hexagon of side one. The obstruction is not a shortage of total
area: the hexagon has the area of six unit triangles. Instead, the same
triangles must cover both the boundary and the interior, and these two
obligations compete.

\readerheading{Seven triangles, seven prescribed points.}
The center and the six vertices of the hexagon are pairwise at distance at
least one. An open unit triangle cannot contain two of them. Thus any
hypothetical cover assigns one triangle to the center and one to each
vertex. We call these the C triangle and the six V triangles. This elementary
observation is the starting point of the proof, not merely a labeling
convention; its precise statement is \zcref{lem:distinct-roles}.

\readerheading{Boundary reach has an interior cost.}
Call a V triangle \emph{supercritical} when the sum of the lengths it covers
on its two incident hexagon edges is greater than one. If the V triangles
cover the entire boundary, neighboring open traces must overlap. Their total
length is therefore greater than six, so at least one V triangle is
supercritical. The crucial local fact is that a supercritical V triangle
misses the midpoint of the segment from its vertex to the center
(\zcref{lem:self-midpoint}). Extra boundary reach creates an interior
obligation for another triangle.

\readerheading{Two branches of the argument.}
If the V triangles cover the entire boundary, two or more supercritical
triangles waste too much area outside the hexagon. Exactly one leads instead
to nine points that the C triangle must contain, but that cannot fit inside
an open unit triangle. If the V triangles leave a boundary gap, the C triangle
must cover it. Reaching the boundary restricts the C triangle to exactly one
radial midpoint. Length budgets then leave only a small number of ways to
cover the other necessary midpoints. These are excluded by five-point or
four-point enclosures, a perimeter estimate, or a replacement that preserves
the boundary and radial segments. The detailed hypotheses are given in
Table~\ref{tab:routing}; this paragraph describes the logic before introducing
the classification notation.

The proof therefore has three geometric measurements: length on a fixed
network, area lost outside the hexagon, and the least equilateral triangle
containing a fixed set of points. The local calculations serve these three
measurements. The exact computer-assisted step is confined to the final
zero-gap enclosure argument; it is not a numerical search over covers.


\subsection{Notation, theorem, and canonical labeling}

All distances, angles, and areas are Euclidean. A unit equilateral triangle
has side length one; its open version is its interior. We write $[x,y]$
for the closed segment, $\conv K$ for the convex hull, and $\interior K$
and $\partial K$ for the interior and boundary. The relative interior
$\operatorname{relint}K$ is its interior within its affine span. The measure $\Haus$ is
one-dimensional Hausdorff measure, which agrees with length on segments.
We use $[t]_+=\max\{t,0\}$ and
$\operatorname{cross}(x,y)=x_1y_2-x_2y_1$ for vectors in $\mathbb R^2$.
A \emph{trace} on a set $E$ means its intersection with $E$.

Let
\[
 O=(0,0),\qquad
 V_i=\left(\cos\frac{i\pi}{3},\sin\frac{i\pi}{3}\right),\qquad
 H=\conv\{V_0,\ldots,V_5\},
\]
where indices lie in $\mathbb Z/6\mathbb Z$.  Write
\[
 e_{i,i+1}=[V_i,V_{i+1}],\qquad r_i=[O,V_i],\qquad M_i=\frac12V_i,
\]
and put
\[
 S=\partial H\cup\bigcup_{i=0}^5r_i.
\]
We call $S$ the \emph{hexagon skeleton}. For $X\in\mathbb R^2$ and $r\ge0$,
write $\mathbb B(X,r)=\{Y:\|Y-X\|\le r\}$ for the closed disk.
\begin{figure}[H]
\centering
\begin{tikzpicture}[scale=2.05]
  \HexCoords
  \DrawHexagon
  \DrawRays
  \fill[CenterBlue] (O) circle (.025);
  \node[figlabel,below left] at (O) {$O$; center role $U_C$};
  \foreach \i in {0,...,5}{
    \fill (V\i) circle (.022);
    \fill[white,draw=RadialTeal,line width=.5pt]
      ($(O)!.5!(V\i)$) circle (.019);
  }
  \node[figlabel,right] at (V0) {$V_0$; role $U_0$};
  \node[figlabel,above right] at (V1) {$V_1$; role $U_1$};
  \node[figlabel,above left] at (V2) {$V_2$; role $U_2$};
  \node[figlabel,left] at (V3) {$V_3$; role $U_3$};
  \node[figlabel,below left] at (V4) {$V_4$; role $U_4$};
  \node[figlabel,below right] at (V5) {$V_5$; role $U_5$};
  \node[figlabel,above] at ($(O)!.5!(V0)$) {$M_0$};
  \node[figlabel,above right] at ($(O)!.68!(V1)$) {$r_1$};
  \node[figlabel,above right] at ($(V0)!.53!(V1)$) {$e_{0,1}$};
\end{tikzpicture}

\caption{The center, vertices, radial segments, and radial midpoints of $H$.}
\label{fig:geometry-roles}
\end{figure}

For a nonempty compact set $K\subset\mathbb R^2$, define its
\emph{equilateral enclosure side} by
\begin{equation}
 \Lambda(K):=\inf\{s>0:K\subset T\text{ for a closed equilateral
 triangle }T\text{ of side }s\}.
 \label{eq:enclosure-definition}
\end{equation}
Then
\begin{equation}
 \Haus(\partial H)=6,\qquad
 \Haus\left(\bigcup_{i=0}^5r_i\right)=6,\qquad \Haus(S)=12.
\label{eq:target-lengths}
\end{equation}

\begin{theorem}[Main theorem]
\label{thm:main}
The hexagon $H$ cannot be covered by seven open equilateral triangles of side
length one.
\end{theorem}

\begin{corollary}[Soifer's Hexagon Problem]
\label{cor:expanded-closed}
For every $L>1$, the regular hexagon $H_L=LH$ cannot be covered by seven closed
unit equilateral triangles.
\end{corollary}

By \zcref{lem:distinct-roles}, a hypothetical cover has a unique
labeling
\[
 U_C,U_0,\ldots,U_5
\]
such that
\[
 O\in U_C,
 \qquad
 V_i\in U_i.
\]
We call $U_C$ the \emph{C triangle} and $U_i$ the \emph{V triangle at
$V_i$}.  Set
\[
 T_C=\overline{U_C},\qquad T_i=\overline{U_i}.
\]
A \emph{role} means one of these labeled triangles. Statements about an
arrangement of roles assume $O\in U_C$ and $V_i\in U_i$; coverage of
$H$, $S$, or $\partial H$ is an additional hypothesis, stated when needed.
A \emph{skeleton cover} covers $S$, not necessarily all of $H$.
An edge is \emph{center-free} when it is disjoint from $U_C$.

\readerheading{Which part of the hexagon must remain covered?}
A full cover also covers the skeleton $S$. The converse is not assumed.
We retain this weaker target because one later replacement preserves the
skeleton, but need not preserve the interior of $H$. Statements about a
skeleton cover are therefore needed for the logic of the replacement, not
just for estimating lengths. The open sets $U_i$ determine coverage;
their closures $T_i$ are used for maxima, incidence classifications, and
upper bounds. A point of $T_i\setminus U_i$ is not covered by $U_i$.




\subsection{How to read the proof}
Section~\ref{sec:structure-common} introduces the incidence classes, actual
reaches, gaps, and midpoint restrictions, in that order. Its final routing
table gives the exhaustive split after the relevant notation is available.
Sections~\ref{sec:trace-length} and \ref{sec:area-loss} establish the length
and area budgets. Section~\ref{sec:finite-enclosure} explains how to fix
witnesses before testing an arbitrary enclosing triangle; the selected-gap
construction is worked through point by point. Section~\ref{sec:assembly}
assembles the obstructions.

The main text states the geometric inputs and explains their use.
Appendices A--E retain the complete local calculations, including both
replacement charts. Appendix F contains the exact tangent-residual
certificate and identifies its accompanying finite data and replay programs.
Reading a calculation can be postponed, but none is being replaced by an
illustration or an experimental assertion.
